\documentclass[11pt]{article}

\usepackage[a4paper,margin=1in]{geometry}
\usepackage{amsmath,amssymb,amsthm}
\usepackage{enumitem}
\usepackage{microtype}
\usepackage{parskip}

% -------------------------------------------------------------------
% Environments
% -------------------------------------------------------------------
\theoremstyle{definition}
\newtheorem*{definition}{Definition}
\newtheorem*{example}{Example}

% -------------------------------------------------------------------
% Notation
% -------------------------------------------------------------------
% A problem cell: \cell{<number>}{<title>}{<points>}
\newcommand{\cell}[3]{%
  \section*{Cell #1: #2\hfill{\normalsize\mdseries #3}}}

\title{Uphill Paths on the Hypercube\\[0.4em]\large Problem Description}
\author{}
\date{}

\begin{document}
\maketitle

% ===================================================================
\section*{Overview}
% ===================================================================

Label the vertices of the hypercube to make as few uphill paths as
possible. IMO 2022 settled the grid; the cube is still open.

\begin{definition}[Labelling]
Let $G$ be a finite simple graph on $n$ vertices. A \emph{labelling}
of $G$ is a bijection
\[
  f \colon V(G) \longrightarrow \{1, 2, \ldots, n\}.
\]
\end{definition}

Fix a labelling.

\begin{definition}[Valleys and uphill paths]
\leavevmode
\begin{itemize}[leftmargin=*]
  \item A vertex $v$ is a \emph{valley} if every neighbour $w$ of $v$
    has
    \[
      f(w) > f(v).
    \]
    An isolated vertex counts as a valley.
  \item An \emph{uphill path} is a sequence
    \[
      (v_1, v_2, \ldots, v_k)
    \]
    with $k \geq 1$ in which $v_1$ is a valley, $v_i$ and $v_{i+1}$
    are adjacent for each $i$, and
    \[
      f(v_1) < f(v_2) < \cdots < f(v_k).
    \]
    A valley on its own ($k = 1$) counts as an uphill path.
\end{itemize}
\end{definition}

\begin{definition}[$U(G)$]
Let $U(G)$ be the smallest possible number of uphill paths, taken
over all labellings of $G$.
\end{definition}

\begin{example}
IMO 2022 Problem~6 asks for $U$ of the $n \times n$ grid graph. The
answer is
\[
  2n^2 - 2n + 1.
\]
\end{example}

\begin{definition}[Hypercube]
This column asks about the $d$-dimensional hypercube $Q_d$. Its
vertex set is
\[
  \{0, 1\}^d,
\]
and two vertices are adjacent when they differ in exactly one
coordinate. $Q_d$ has $2^d$ vertices and $d \cdot 2^{d-1}$ edges.
\end{definition}

% ===================================================================
\section*{What to Hand In}
% ===================================================================

\begin{itemize}[leftmargin=*]
  \item \textbf{C1 to C4.}
    Hand in the values, and for each value an explicit labelling that
    attains it as a list of the $2^d$ vertices in increasing label
    order, written as $0/1$ strings.

    These cells are marked correct on the values alone, but a value
    with no labelling behind it will not survive the later cells,
    which build on the construction.
  \item \textbf{C5.}
    Hand in either a labelling of $Q_9$ in the same format, whose
    uphill paths will be counted mechanically, or a proof of the
    lower bound.
  \item \textbf{C6.}
    Hand in all three of:
    \begin{enumerate}[leftmargin=*]
      \item the value;
      \item an explicit labelling that attains it, in the same format;
      \item a proof that no labelling gives fewer uphill paths.
    \end{enumerate}
  \item \textbf{Computation.}
    A lower-bound proof may be computer-assisted. If it is, include
    the code; it must run in under 10 minutes on a laptop, and you
    must explain why the computation proves the bound.
  \item \textbf{Status.}
    Say clearly which cells you consider solved and which are
    partial.
\end{itemize}

% ===================================================================
\cell{1}{$U(Q_3)$ and $U(Q_4)$}{1 point}
% ===================================================================

\textbf{Checked instantly.}

The two smallest interesting cubes: $Q_3$ has $8$ vertices and $12$
edges, $Q_4$ has $16$ vertices and $32$ edges.

Determine $U(Q_3)$ and $U(Q_4)$.

% ===================================================================
\cell{2}{$U(Q_5)$}{2 points}
% ===================================================================

\textbf{Checked instantly.}

$Q_5$ has $32$ vertices and $80$ edges. Determine $U(Q_5)$.

% ===================================================================
\cell{3}{$U(Q_6)$}{3 points}
% ===================================================================

\textbf{Checked instantly.}

$Q_6$ has $64$ vertices and $192$ edges. Determine $U(Q_6)$.

% ===================================================================
\cell{4}{$U(Q_7)$ and $U(Q_8)$}{5 points}
% ===================================================================

\textbf{Checked instantly.}

Two larger cubes: $Q_7$ has $128$ vertices and $448$ edges, $Q_8$ has
$256$ vertices and $1024$ edges.

Determine $U(Q_7)$ and $U(Q_8)$.

% ===================================================================
\cell{5}{Bounds for $U(Q_9)$}{8 points}
% ===================================================================

\textbf{Judged.}

$Q_9$ has $512$ vertices and $2304$ edges. The best bounds known to
the organisers are
\[
  2368 \leq U(Q_9) \leq 2400;
\]
the lower bound is unpublished.

Improve either one: prove that
\[
  U(Q_9) \geq 2369,
\]
or exhibit a labelling of $Q_9$ with at most $2399$ uphill paths.

% ===================================================================
\cell{6}{$U(Q_9)$}{13 points}
% ===================================================================

\textbf{Judged. Open question.}

The same cube, settled completely. Determine $U(Q_9)$ exactly, with
proof of both bounds.

\end{document}
